\chapter{Glossary}\label{ch:glossary}
This chapter is a glossary of the technical terms used in every chapter. Each entry gives a short definition and a pointer to the place where the term is defined in full. The symbols for the global and local data are in Chapter~\ref{ch:notation}, and the notation specific to one chapter is listed there as well. The chapter proves nothing and can be consulted at any point.

\section{Covers and the cocycle normalization}\label{gl:cocycle}
This section lists the terms that describe covering groups and fixes the cocycle normalization used in every chapter. The book compares covers built in different ways, and a comparison of constants is only meaningful once one normalization is fixed.

\begin{description}
\item[Cover, $p$, $\overline{H'}$]\index{cover p h@Cover, $p$, $\overline{H'}$} An \emph{$n$-fold cover} of a
locally compact group $H$ is a topological central extension
$1\to\mu_n\to\overline H\xrightarrow{\ p\ }H\to1$.  For a subgroup
$H'\subset H$ we write $\overline{H'}:=p^{-1}(H')$.  Chapters~\ref{ch:I} and~\ref{ch:II}
write $\overline H$; Parts~III, IV and the companion chapter write
$\widetilde H$, $\widetilde Z$.
\item[Commutators]\index{commutators@Commutators} For lifts $\tilde x,\tilde y$ of commuting
elements, $[x,y]:=\tilde x\tilde y\tilde x^{-1}\tilde y^{-1}\in\mu_n$;
it does not depend on the lifts, and it is bimultiplicative on a
torus \cite[\S3]{FK86}.
\item[Kazhdan--Patterson cocycle $\sigma_F$]\index{kazhdan patterson cocycle f@Kazhdan--Patterson cocycle $\sigma_F$} The cocycle of
\cite[\S0.1]{KP84} and \cite[\S2]{FK86} on
$\mathrm{GL}(r,F)$ with twisting parameter $0$.  On the diagonal
torus,
\[
\sigma_F(\mathrm{diag}(a_i),\mathrm{diag}(b_i))=\prod_{i<j}(a_i,b_j)_F.
\]
\item[Parameter-$c'$ cover $\overline G^{(c')}$]\index{parameter c cover g c@Parameter-$c'$ cover $\overline G^{(c')}$} The cover of
$\mathrm{GL}(r,F)$ with cocycle
\[
\sigma_F(g,h)\,(\det g,\det h)_F^{\,c'}.
\]
The comparison uses
$c'=c_v$; the parameter is forced (\cref{I:rem:forced}).  On the
diagonal torus the commutator of $\mathrm{diag}(t_i)$ and
$\mathrm{diag}(u_j)$ is
\[
\prod_{i,j}(t_i,u_j)_F^{A_{ij}},\qquad A=-I_r+(1+2c_v)J_r,
\]
where $J_r$ is the all-ones matrix (\S\ref{IV:subsec:restriction}).
\item[Ambient cover $\overline{\mathcal G}_v$]\index{ambient cover g v@Ambient cover $\overline{\mathcal G}_v$} The cover of
$\mathcal G_v=\mathrm{GL}(r,L_v)=\prod_{w\mid v}\mathrm{GL}(r,L_w)$
with cocycle
\[
\prod_{w\mid v}\sigma_{L_w}(g_w,h_w)\,(\det g_w,\det h_w)_{L_w}^{\,c_v},
\]
the copies of $\mu_n$ being identified (\S\ref{I:subsec:ambient} and
equation~\eqref{I:eq:ambient}).  Both covers compared in the series are subgroups of
it, so their common central subgroup $p^{-1}(K_v^\times)$ is literally
the same subgroup on the two sides.
\item[Restriction and relabeling]\index{restriction and relabeling@Restriction and relabeling} Pulled back along
$\iota_D:D_v^\times\hookrightarrow\mathcal G_v$, the ambient cocycle
is the local component of the cocycle
$\beta=\sigma_L\cdot(\mathrm{Nrd},\mathrm{Nrd})^c$ of the corrected
1988 paper \cite[equation~(4)]{JC26a}.  Pulled back along the diagonal
embedding of $\mathrm{GL}(r,K_v)$, it is isomorphic, \emph{as an
extension}, to the extension with cocycle
$\sigma_F^{\,r}(\det,\det)_F^{\,c}$ (\cref{I:prop:restriction}).  The
group $\mu_n$ is relabeled \emph{once}, at every place and globally,
by $\zeta\mapsto\zeta^{r'}$; after relabeling the second extension is
$\overline G^{(c_v)}$.  Formulas state whether they hold before or
after relabeling.  The isomorphism is one of extensions, not an
equality of cocycles, and it is unique up to
$\tilde g\mapsto\chi(\det p(\tilde g))\,\tilde g$ with
$\chi:F^\times\to\mu_n$ continuous.  Central characters are therefore
never transported through it (\cref{I:rem:chidet}).
\item[Scalar centrality]\index{scalar centrality@Scalar centrality} Because $n\mid r-1+2rc_v$, the whole inverse
image $\overline Z=p^{-1}(F^\times)$ of the scalars is central
(\cref{I:lem:center}).  A cover with this property is called
\emph{scalar-central}.
\item[Global normalization]\index{global normalization@Global normalization} Globally one uses the normalized cocycle
\[\tau_w(g,h)=\sigma_w(g,h)\,\kappa_w(g)\kappa_w(h)\kappa_w(gh)^{-1}\]
of \cite[\S2]{FK86}, where $\kappa_w$ is trivial for $w\mid
n\infty$.  The compact splitting is $\kappa^*_w(k)=(k,1)_\tau$ on
$\mathrm{GL}(r,\mathcal O_{L_w})$ for $w\nmid n\infty$, and the
rational splitting is $\iota(x)=s_\tau(x)\,\kappa_L(x)^{-1}$; in the
unnormalized models the rational splitting is $x\mapsto(x,1)_\sigma$
at every place (\cref{I:lem:splittings}).  The global ambient cover
$\overline{\mathcal G}(\mathbb A_L)$ has cocycle
$\tau_L\cdot(\det,\det)_{\mathbb A_L}^{\,c_v}$.
\item[Star map]\index{star map@Star map} For $h$ in the base group and any lift $\tilde h$,
$h^*:=\tilde h^{\,n}$.  It does not depend on the lift, it is
conjugation equivariant, $(ah)^*=a^*h^*$ for scalars $a$, and it is
multiplicative on every torus because $n$ is odd (\cref{III:def:star}; \S\ref{IV:subsec:star}).  For $\zeta\in\mu_n(F)$ one has
$\zeta^*=1$.
\end{description}

\section{Genuine objects, central characters and good elements}\label{gl:genuine}
This section lists the terms for objects on a cover: genuine representations and functions, their central characters, and good elements. Harmonic analysis on a cover takes place on the good elements, so these terms are used in almost every statement of the book.

\begin{description}
\item[Genuine representation]\index{genuine representation@Genuine representation} A representation of $\overline H$ on
which $\mu_n$ acts by the tautological character $\zeta\mapsto\zeta$.
\item[$\tilde\theta$, $\theta$]\index{$\tilde\theta$, $\theta$@$\tilde\theta$, $\theta$} A genuine character $\tilde\theta$ of
$\overline Z$ (unitary from \S\ref{I:subsec:globalcovers} on; globally, a character of
$\overline Z(\mathbb A_K)$ trivial on $\iota(K^\times)$).  The same
$\tilde\theta$ serves both covers, since $\overline Z$ is a common
subgroup.  The ordinary character attached to it is
$\theta(a):=\tilde\theta(a^*)$; it is trivial on $\mu_n$, and
$\theta$ does not determine $\tilde\theta$ (\S\ref{IV:subsec:star}).
Chapter~\ref{ch:III} requires only that $\tilde\theta|_{\mu_n}$ be faithful.
\item[Genuine test functions]\index{genuine test functions@Genuine test functions} $C_c^\infty(\overline H,\tilde\theta)_{\mathrm{gen}}$
is the space of smooth functions with support compact modulo
$\overline Z$ and $f(z\tilde h)=\tilde\theta(z)^{-1}f(\tilde h)$ for
$z\in\overline Z$; in particular $f(\zeta\tilde h)=\zeta^{-1}f(\tilde h)$.
Characters and invariant class functions of genuine representations
transform by $\tilde\theta$ itself (\S\ref{III:sec:setting}).
\item[Regular, elliptic]\index{regular elliptic@Regular, elliptic} An element of $\mathrm{GL}(r,F)$ is
\emph{regular} if its eigenvalues are distinct, and \emph{elliptic
regular} if moreover its characteristic polynomial is irreducible.
Since $r$ is prime, the regular elements of $D_v^\times$ are the
noncentral ones, and they are all elliptic.  An element of a cover is
regular (elliptic) if its image is.
\item[Good element]\index{good element@Good element} A regular $\tilde\gamma$ is \emph{good} if it
commutes with $p^{-1}(T_\gamma)$, where $T_\gamma$ is the centralizer
of $\gamma=p(\tilde\gamma)$; equivalently, $\tilde\gamma$ is central
in the covered torus $\overline{T_\gamma}$.  Genuine invariant
functions vanish at regular elements that are not good.  For a
maximal subfield $E$ of $D_v$ (or an elliptic torus $E^\times$ of
$\mathrm{GL}(r,F)$), the good elements of $\overline{E^\times}$ are
$p^{-1}(F^\times E^{\times n})=\overline Z\cdot\{t^*\}$, of index
$|n|_F^{-(r-1)}$, and every genuine irreducible representation of
$\overline{E^\times}$ has dimension $|n|_F^{-(r-1)/2}$ (\cref{I:prop:good}).  Covered elliptic tori are abelian exactly at the
places not dividing $n$.
\end{description}

\section{Orbital integrals and matching}\label{gl:matching}
This section lists the terms used to compare test functions on the two groups: the normalized orbital integrals, the identification of tori, matching, and germs. Matching of orbital integrals is the local input to the comparison of trace formulas.

\begin{description}
\item[$\Delta$]\index{$\Delta$@$\Delta$} For regular semisimple $\gamma$,
$\Delta(\gamma)=|D(\gamma)|_F^{1/2}|\det\gamma|_F^{-(r-1)/2}$, where
$D(\gamma)$ is the discriminant of the (reduced) characteristic
polynomial and $\det=\mathrm{Nrd}$ on the division-algebra side.
Thus $\Delta(\gamma)^2$ is the Weyl discriminant
$|\det(1-\mathrm{Ad}(\gamma^{-1}))|_{\mathfrak g/\mathfrak t}|_F$
(\cref{I:def:oi} and \cref{I:lem:traceOI}).
\item[$\Phi$, $F=\Delta\Phi$]\index{f@$\Phi$, $F=\Delta\Phi$} $\Phi(\tilde\gamma,\tilde f)$ is the
orbital integral over $T_\gamma\backslash G$ of a genuine test
function at a good regular $\tilde\gamma$, and
$F(\tilde\gamma,\tilde f)=\Delta(\gamma)\Phi(\tilde\gamma,\tilde f)$
is the normalized orbital integral, set to $0$ at regular elements
that are not good (\cref{I:def:oi}).  Chapter~\ref{ch:III} writes
$J(t,f)=\Delta(t)\,O(t,f)$ for the same function, with its own
measure conventions.
\item[Canonical identification $\jmath$]\index{canonical identification@Canonical identification $\jmath$} At $v\in S$ and for a maximal
subfield $E\subset D_v$ with an embedding $j:E\hookrightarrow
M_r(F)$, every isomorphism of covered tori that is the identity on
$\overline Z$ has the same restriction $\jmath_{E,j}$ to the good
elements; these restrictions induce a bijection $\jmath$ from the good
regular classes of the cover of $D_v^\times$ to the good regular
elliptic classes of the cover of $\mathrm{GL}(r,F)$ (\cref{I:thm:ident} and \cref{I:cor:classes}).
\item[Matching]\index{matching@Matching} Genuine test functions $\tilde f_D$ and $\tilde f$
\emph{match} if (M1)
$F(\tilde\gamma,\tilde f_D)=(-1)^{r-1}F(\jmath(\tilde\gamma),\tilde f)$
at every good regular $\tilde\gamma$, and (M2) $F(\tilde\delta,\tilde
f)=0$ at every good regular non-elliptic $\tilde\delta$ (\cref{I:def:matching}).  The sign sits in the matching of functions, so that
the expected character relation is
$\Theta_{\overline\pi_D}=(-1)^{r-1}\Theta_{\overline\pi}\circ\jmath$.
At a place $v\notin S$, matching is simply the identity map in the
coordinates fixed by $\psi_v$, with no sign and nothing to solve; see
Example~\ref{I:ex:split-matching} for this worked out at a named
split place of the running example.
\item[Supercusp form]\index{supercusp form@Supercusp form} A genuine test function on the cover of
$\mathrm{GL}(r,F)$ whose integrals along the unipotent radicals of
all proper parabolic subgroups vanish, for example a matrix
coefficient of a supercuspidal representation \cite[\S13]{FK86}.
\item[Germ]\index{germ@Germ} The asymptotic behavior of orbital integrals near a
singular element, expressed through Shalika germs \cite[\S8]{FK86}; used near the center in Chapter~\ref{ch:II}.
\end{description}

\section{Haar measures}\label{gl:measures}
This section records the Haar measures used in each chapter. The chapters normalize measures differently, and constants such as formal degrees depend on the choice, so the conventions are listed side by side.

\begin{description}
\item[Chapter~\ref{ch:I}]\index{chapter ch i@Chapter~\ref{ch:I}} At $v\in S$ the Haar measures on $D_v^\times/F^\times$
and $\mathrm{GL}(r,F)/F^\times$ are arbitrary but fixed, and matching
is defined with respect to them; torus measures are transported by
conjugation and do not affect matching.  At $v\notin S$ the measures
correspond under the identifications $\psi_v$ (\S\S\ref{I:subsec:oi}, \ref{I:subsec:measures}).
\item[Chapter~\ref{ch:II}]\index{chapter ch ii@Chapter~\ref{ch:II}} At $v\in S$,
$\mathrm{vol}(D_v^\times/F^\times)=d(\mathrm{St})^{-1}$, where
$d(\mathrm{St})$ is the formal degree of the Steinberg representation
for the chosen measure on $\mathrm{GL}(r,F)/F^\times$ (\cref{II:def:measures}).  Example~\ref{II:ex:trivgerm} computes the resulting
trivial-germ sign at the running example's tame place.
\item[Chapter~\ref{ch:III}]\index{chapter ch iii@Chapter~\ref{ch:III}} For the transform $L_D$, the quotients $H/Z$ and
$T/Z$ carry their Haar \emph{probability} measures (\S\ref{III:sec:local}).
\item[Formal degree]\index{formal degree@Formal degree} For a representation $\sigma$ that is square
integrable modulo the center, the constant $d(\sigma)$ in the Schur
orthogonality relations; it depends on the Haar measure.
\end{description}

\section{The division-side transform}\label{gl:transform}
This section lists the terms attached to the transform $L_D$ of class functions on the division-algebra side. The transform is the main object of Chapters~\ref{ch:III}--\ref{ch:W}.

\begin{description}
\item[$H$, $\widetilde H$, $L_D$]\index{h h l d@$H$, $\widetilde H$, $L_D$} For a local division algebra $D$ of
degree $r$ over $F$, $H=D^\times$ and $\widetilde H$ is its cover from
Chapter~\ref{ch:I}.  The transform $L_D$ sends $\theta$-equivariant invariant
$L^2$ functions $\varphi$ on $H$ to genuine
$\tilde\theta$-equivariant invariant functions on $\widetilde H$ by
$(L_D\varphi)(z\,t^*)=b_D\,\tilde\theta(z)\,\Delta(t)\Delta(t^n)^{-1}\varphi(t)$
at good regular elements, and $0$ at the others; it is unitary
(\cref{III:def:L} and \cref{III:thm:unitary}).
\item[$b_D$, $c_D$]\index{b d c d@$b_D$, $c_D$} $b_D=|n|_F^{-(r-1)/2}$ and
$c_D=|n|_F^{-(r^2-1)/2}$ (\S\S\ref{III:sec:local}--\ref{III:sec:characters}).  Both equal $1$ at a
tame place; in the running example $b_D=3$ and $c_D=27$ at
$\mathfrak p_3$.
\item[Integrality of $L_D$]\index{integrality of l d@Integrality of $L_D$} The question whether $L_D$ carries
irreducible characters to irreducible genuine characters.  It holds at
tame places (Chapter~\ref{ch:III}) and is open in general at wild places
(Chapter~\ref{ch:IV}; companion chapter).
\end{description}

\section{Named results and hypotheses}\label{gl:named}
This section lists the theorems and hypotheses that the book refers to by name, with the place where each is stated. It lets a reader who meets a name such as Theorem~C or $(\mathrm{Hyp}_D)$ find its statement quickly.

\begin{description}
\item[Theorems A, B, C]\index{theorems a b c@Theorems A, B, C} Chapter~\ref{ch:I}: the comparison of covered elliptic
tori, the matching of regular-supported test functions, and the
global transfer under $(\mathrm{Hyp}_D)$ (with its converse under (Hyp)).
\item[Theorems D, E, C$'$]\index{theorems d e c@Theorems D, E, C$'$} Chapter~\ref{ch:II}: the equivariant section near the
center and its compatibility with $\jmath$; the matching of arbitrary
genuine test functions; and the global transfer under the weaker
hypothesis $(\mathrm{Hyp}_D')$ (with its converse under $(\mathrm{Hyp}')$).
\item[$(\mathrm{Hyp}_D)$, (Hyp), $(\mathrm{Hyp}_D')$, $(\mathrm{Hyp}')$]\index{hyp d hyp hyp d hyp@$(\mathrm{Hyp}_D)$, (Hyp), $(\mathrm{Hyp}_D')$, $(\mathrm{Hyp}')$}
Non-cancellation hypotheses: the character of a near-equivalence
packet at the places of $S$ is not identically zero on the relevant
matching test functions (\cref{I:def:HD} and
\cref{I:thm:converse}; \cref{II:def:HDprime} and \S\ref{II:subsec:Cprime}).  $(\mathrm{Hyp}_D)$
holds when the $D$-side packet is finite (\cref{I:prop:HDfinite}); sufficient conditions for (Hyp) are in \cref{I:rem:H}; and $(\mathrm{Hyp}_D)\Rightarrow(\mathrm{Hyp}_D')$,
$(\mathrm{Hyp})\Rightarrow(\mathrm{Hyp}')$ (\S\ref{II:subsec:Cprime}).
\item[(D1)--(D4)]\index{d1 d4@(D1)--(D4)} The steps left open by Chapter~\ref{ch:I}: matching of
arbitrary functions, local character relations, finiteness of
near-equivalence packets, and composition with \cite{JC26a} (\S\ref{I:sec:deferred}).
Their status in the later Parts is recorded in \S\ref{I:subsec:partII}.
\item[Toral positivity; the discriminant ratio]\index{toral positivity@Toral positivity}\index{discriminant ratio@Discriminant ratio $\Delta(t)/\Delta(t^n)$} For a regular
element $t$ in a maximal field $E\subset D_v$, write
$\Delta(t)=\lvert\det(1-\mathrm{Ad}(t)\mid\mathfrak h/\mathfrak
t)\rvert_F^{1/2}$ for the Weyl-type discriminant factor and
$R_E(t):=\Delta(t)/\Delta(t^n)$ for the \emph{discriminant ratio}
(\eqref{IV:eq:Delta}); this ratio is locally constant on the compact
torus $A_E=E^\times/F^\times$ and so has a finite Fourier expansion
there. \emph{Toral positivity} is the hypothesis that this expansion
has only nonnegative integer coefficients, equivalently that $R_E$ is
the character of a finite permutation representation of $A_E$;
by Proposition~\ref{IV:prop:int} this is equivalent to integrality of
the entire unitary character-transform matrix $(a_{\rho\tau})$ of
$L_D$ (\S\ref{gl:transform} below), which is what Chapters~\ref{ch:III}
and~\ref{ch:IV} need but cannot prove in general. Chapter~\ref{ch:IV}'s
global comparison is conditional on toral positivity holding at every
place it is needed; Chapter~\ref{ch:W} proves it unconditionally for
every quadratic $E$ at the running example's wild place
(Theorem~\ref{W:thm:main}), with explicit multiplicities $(3,2,2,2)$
on the unramified torus and $(3,2,1,1,1,1)$ on either ramified one.
\item[Targets A, B, C]\index{targets a b c@Targets A, B, C} The companion chapter's three local problems at
the wild place of the running example and beyond (Chapter~\ref{ch:W}, \S\ref{W:sec:intro}).
\end{description}

